% conclusion.tex (chundawat-h)

\section{Conclusion}
\label{sec:conclusion}

We compared a rule-based baseline, a text-only Transformer (BERT) and a
layout-aware Transformer (LayoutLM) for extracting company, date,
address and total from scanned receipts. Rules already solve dates
(95.8\% $F_1$) but struggle with fields whose position and wording vary,
reaching 74.8\% macro $F_1$. A fine-tuned text-only BERT raises this to
91.9\%, and adding the 2D word positions (LayoutLM) to 93.7\%.

So, to answer our research question, layout adds a small but
statistically reliable 1.8 points over text alone. The gain is
concentrated on company ($+3.6$) and address ($+2.6$), whose extent on
the page is defined visually. It is absent for the total, where the
training labels do not consistently mark the amount on the ``TOTAL''
line. Most of the improvement over rules comes from contextual text
features; layout mainly fixes block-boundary errors. Our two-gold
evaluation further shows that OCR noise, not the tagging model, is the
main limit on company and address: even perfect tagging reaches only
80.4\% and 58.2\% against the official answers.

\paragraph{Limitations and future work.} Each model was trained once
(seed 42); the bootstrap covers test-set sampling but not training
randomness, so several seeds would make the per-field gains more
reliable. Natural next steps are re-labelling the total on the
``TOTAL'' line to give layout a consistent signal, models that also use
the image (LayoutLMv2~\cite{xu2021layoutlmv2}), OCR post-correction to
raise the ceiling, and evaluation on other receipt datasets such as
CORD~\cite{park2019cord}.
